%%%%%%%%%%%%%%%%%%%%%%%%%%%%%%%%%%%%%%%%%%%%%%%%%%%%%%%%%%%%%%%%%%%%%%%%%%%%%%%%
% 架构重建与受控变异：面向先进大语言模型安全防御的统一框架
% 中文版 | Chinese Version
% 格式：USENIX (usenix-2020-09.sty)
%%%%%%%%%%%%%%%%%%%%%%%%%%%%%%%%%%%%%%%%%%%%%%%%%%%%%%%%%%%%%%%%%%%%%%%%%%%%%%%%

\documentclass[letterpaper,twocolumn,10pt]{article}
\usepackage{CJKutf8}
\usepackage{usenix-2020-09}

\usepackage{amsmath,amssymb,amsthm}
\usepackage{booktabs}
\usepackage{multirow}
\usepackage{xspace}
\usepackage{graphicx}
\usepackage{float}

%-------------------------------------------------------------------------------
\title{\Large \bf 架构重建与受控变异：面向先进大语言\\
  模型安全防御的统一框架}

\author{
{\rm Author One}\\
Anthropic \and
{\rm Author Two}\\
Google
\thanks{以跨学科教授及 Anthropic/Google 首席安全架构师视角撰写；作者姓名可按需替换。}
}

\date{}

%-------------------------------------------------------------------------------
\begin{document}
\begin{CJK}{UTF8}{gbsn}
%-------------------------------------------------------------------------------
\maketitle

%-------------------------------------------------------------------------------
\begin{abstract}
大语言模型（LLM）已从研究工具转变为关键系统核心组件，传统安全范式随之承压。本文提出一个统一框架，将\emph{架构重建}（通过自改进设计、上下文构建防御与混合裁判集成的运行时自适应）与\emph{受控变异}（以元认知监控与示例引导的银标流水线为锚）相结合。我们形式化二者之间的互动关系，并将安全与脆弱的 LLM 自我反思\emph{解耦}：自适应由上下文混淆（诱饵、ID 对齐）、黄金示例与经过校准的裁判网络驱动，而非由模型「自我左脚踩右脚」。我们将 SISF、上下文构建防御（ID 对齐、诱饵注入、末端自我提醒）、带分布评估与强制人工校准的混合裁判集成、价值/本体锚定、元认知监控与示例引导变异，以及槽位变异与银标数据流水线，综合为单一、面向敌对审查加固的架构。我们讨论威胁模型（含任务重构、心理操纵与 AI 裁判偏差）、指标与治理，并论证：要实现以机器速度运行的、可扩展且可审计的 LLM 安全，此类深度重构不可或缺。
\end{abstract}

%-------------------------------------------------------------------------------
\section{引言}
%-------------------------------------------------------------------------------

大语言模型已不再是边缘研究产物，而是关键基础设施的支撑：代码生成、具系统访问权的智能体、谈判与合规敏感工作流。这一转变暴露出根本性错配：传统软件工程与安全保证假定组件是确定性的、可穷举规约的，而 LLM 驱动系统呈现 emergent、随机行为，无法在设计时完全规定~\cite{sisf-arxiv,agentos-arxiv}。静态防御与周期性人工红队以人类速度运转，难以跟上自动化对抗生成；一旦新型越狱或提示注入得逞，架构往往保持不变，仍对同类攻击脆弱。

两个概念支柱应对这一缺口。\textbf{架构重建}指系统安全层的自主、运行时自适应——将安全从一次性部署前审计转变为持续的、由 AI 驱动的过程，在检测到失败时合成并部署新策略~\cite{sisf-arxiv}。\textbf{受控变异}指对提示、示范与基准进行严格、语义保持的扰动，以诊断鲁棒性、将专家行为迁移到智能体并扩展评测覆盖，同时避免测试集污染~\cite{reflective-mutation,adm-es-arxiv}。二者单独均不足：缺乏有原则的变异基质的重建易对狭窄威胁模型过拟合；缺乏架构自适应的变异则在每轮评测后仍使防御处于静态。

\textbf{安全解耦。}本框架明确承认：在对抗压力（任务重构、心理操纵、裁判偏差）下，LLM 的认知与反思是脆弱的。因此我们将安全与模型无约束的自我改进\emph{解耦}。自适应不是 LLM 盲目「左脚踩右脚」，而是通过\emph{上下文混淆}（诱饵与 ID 对齐）剥夺攻击者的清晰操作空间，通过\emph{黄金示例}锚定防守方价值，并通过带强制人工校准的\emph{多样化裁判集成}保证进化方向正确。

本文将二者统一于单一框架下，并采用跨学科学界与产业首席安全架构师（Anthropic 与 Google）的双重角色以衔接理论与部署。主要贡献如下：

\begin{itemize}
\item \emph{统一概念模型}：将架构重建与受控变异联系起来，明确边界（资产、Oracle、反馈环）与威胁模型（\S\ref{sec:framework}）。
\item \emph{重建轴}：SISF、上下文构建防御（ID 对齐、诱饵注入、末端自我提醒）、带分布评估与强制人工校准的混合裁判集成，以及价值/本体锚定（\S\ref{sec:reconstruction}）。
\item \emph{变异轴}：元认知监控与对抗性重置、示例引导变异（黄金数据集与分层经验记忆）、槽位变异、银标流水线（ADM-ES）及元数据层控制（\S\ref{sec:mutation}）。
\item \emph{整合、指标与治理}：重建与变异如何相互反馈、评测协议及对 attestation 与发布门控的意涵（\S\ref{sec:integration}）。
\end{itemize}

%-------------------------------------------------------------------------------
\section{统一框架与威胁模型}
\label{sec:framework}
%-------------------------------------------------------------------------------

\subsection{资产与 Oracle}

\emph{资产}是系统有效的「安全表面」：策略、护栏、触发集以及用于评估与调优它们的数据。\emph{Oracle}是对行为进行标注（安全/不安全、合规/不合规）的过程。在多轮设定下，评估者（或对手）与系统交互并观察结果——二值决策、分数或失败报告——形成统计查询接口。若同一项目在多轮中不受控地暴露，会助长重建攻击，恢复触发模式或基准内容~\cite{llm-reconstruction-emergent,ControlledMutation_SafetyEval}。

\subsection{重建与变异作为双杠杆}

\textbf{架构重建}改变\emph{谁}在\emph{何时}执行安全：从固定的人写规则变为在运行时更新的自适应模块（裁判、策略合成、反馈环）。其本身不控制暴露给 Oracle \emph{什么}数据。\textbf{受控变异}改变\emph{暴露什么}：系统不再重复提供规范项目，而是由变异器提供语义等价变体，并辅以轮次调度与分散，以阻止基于梯度或优化的逆推~\cite{ControlledMutation_SafetyEval}。二者结合：重建提高自适应上限，变异保护资产并保持 Oracle 信息量。图~\ref{fig:framework} 给出统一架构概览。

\subsection{威胁模型与敌对审查攻击面}

我们考虑（i）\emph{外部}威胁：安全数据的多轮抽取、越狱与提示注入外泄；（ii）\emph{内部}威胁：合成污染、自动标注或变异数据中的偏差，以及认知操纵~\cite{epistemic-grounding}。对\emph{敌对审查}红队，三个攻击面需要深层次架构（而非仅参数）响应：

\textbf{任务重构：}对手将抽取伪装成合法上下文处理任务，使基于意图的过滤器只看到「正常」行为。防御必须污染攻击者操作空间（上下文构建防御，\S\ref{sec:reconstruction}）。

\textbf{心理操纵（HPM）：}煤气灯与认知压迫攻击诱发模型自我怀疑与价值漂移；反思与自我诊断进而可产生\emph{恶意}变体。防御必须将变异锚定于外部认知锚点与独立监控（\S\ref{sec:mutation}）。

\textbf{AI 裁判偏差：}单一动态裁判存在系统性盲区；反馈可能将闭环导向错误方向。防御须采用多样化裁判集成与强制人工校准（\S\ref{sec:judge}）。防御须保持\emph{评测保真度}与\emph{可审计性}。

\begin{figure*}[t]
\centering
\includegraphics[width=\textwidth]{Arch}
\caption{\textbf{统一框架概览。}架构将\emph{架构重建}（左：SISF、上下文构建防御、带价值锚定的混合裁判集成、策略合成与反馈环）与\emph{受控变异}（右：元认知监控、带黄金数据集与经验记忆的示例引导变异、轮次调度）耦合。资产与 Oracle 位于边界；重建决定谁在何时执行安全，变异控制暴露什么。强制人工校准与异质恢复通道闭合裁判偏差与防御性降级的闭环。}
\label{fig:framework}
\end{figure*}

%-------------------------------------------------------------------------------
\section{架构重建}
\label{sec:reconstruction}
%-------------------------------------------------------------------------------

\subsection{自改进安全框架（SISF）}

SISF 将安全重述为持续过程~\cite{sisf-arxiv}。未对齐基座 LLM 提供能力；违规检测不交给单一裁判，而交给\emph{混合裁判集成}（\S\ref{sec:judge}）；策略合成模块消费共识信号并生成新的启发式与语义策略；动态反馈环即时部署。因此\emph{架构}随失败而自我重建。在 AdvBench 式评估中，当裁判流水线经适当校准时，此类框架可将攻击成功率从 100\% 降至约 45\%，同时在良性输入上保持 0\% 假阳性率~\cite{sisf-arxiv}。要点是：当新攻击得逞时，系统并非保持静态，而是在同一运行环内合成并应用新防御。

\subsection{上下文构建防御（超越意图审查）}

当对手采用\emph{任务重构}时，基于意图的「语义防火墙」失效：抽取被伪装成合法上下文处理任务，恶意意图不再可被检测~\cite{agentos-arxiv}。我们升级为\textbf{上下文构建防御}，不仅过滤输入，而且\emph{主动污染攻击者的操作空间}。

\textbf{ID 对齐：}在检索到的上下文或中间结果到达未对齐基座 LLM 之前，\textbf{ID 对齐}模块剥离或规范化所有实例标识符。攻击者无法再通过指定 ID 精确定界并抽取敏感片段；良性问答不受影响。在不改变名义任务语义的前提下降低重建精度。

\textbf{动态诱饵注入与风险评分：}系统注入带\textbf{密码学水印}（如基于哈希的结构）的\emph{诱饵字段}（合成实体或关联），使偶然幻觉复现概率极低。后端指令将诱饵标为内部禁止。良性用户忽略它们；任务重构型攻击者常覆盖指令并暴露诱饵数据。为避免不可接受的 FPR，采用\textbf{诱饵命中密度}与\textbf{风险评分}：单次低保真触碰触发警告与意图重校验；仅当对完整诱饵结构的高保真抽取才触发会话终止与审计。我们不声称密码学隔离或残余抽取「低于实际效用」：实证工作（如 GRASP 式评估）表明，即便有 ID 对齐、诱饵与拒绝门控，\textbf{残余抽取}（如 F1 约 15\%）仍存在~\cite{subgraph-rag-attack}；在 PBU 下复合防御可基本失效~\cite{pbu-reconstruct}。设计目标是\textbf{指数级提高攻击者查询复杂度}：诱饵与 ID 对齐使达到给定重建保真度所需轮次增长（实践中常超线性）。诱饵命中密度对\emph{暴力}抽取有效；但我们不断言其能在具有经济价值的重建完成前可靠触发熔断。PBU 攻击者\textbf{完美模仿良性用户分布}：其诱饵触碰率可与良性幻觉的 FPR 相当——阈值过低导致不可接受的假阳性，过高则允许慢速拼接至 0.75+ 语义泄露。因此我们承认在面对完美 PBU 时\textbf{可用性与安全性之间的不可避免权衡}。框架将其重述为\textbf{经济防御}：目标是在资产可被轮换或弃用的前提下，使慢速 PBU 抽取的\emph{成本}（算力、Token、时间）超过被保护资产的\textbf{有效生命周期}。我们承认这\textbf{依赖资产可轮换性假设}。对\textbf{静态、长期或永久}高价值资产（如核心 IP、患者 PII 或基础对齐数据集），基于生命周期的经济学不适用——资产无法在任意有意义的时间窗内被「轮换」。对此类资产，系统不能仅靠经济学作为终极屏障，而必须在元数据与访问控制层（\S\ref{sec:metadata}）实施\textbf{基于角色的访问控制（RBAC）}与\textbf{信息流追踪}，将经济学视为\textbf{纵深防御}的一层而非唯一机制。

\textbf{末端自我提醒：}在多轮交互的最后一轮，将安全指令强制附加到用户查询作为系统提示，在生成时强化边界意识并减少会话末漂移。

\subsection{混合裁判集成与校准}
\label{sec:judge}

单一 AI 裁判存在系统性偏差与盲区。我们以\textbf{混合裁判集成}替代：多个不同架构与训练分布的裁判（如 Llama-Guard~\cite{llama-guard}、StrongREJECT~\cite{strongreject} 或任务专用分类器）。仅当裁判达成\emph{共识}时，反馈信号才传入策略合成模块。我们不假定集成多样性单独消除偏差：近期工作表明\textbf{格式敏感性}（空格、系统提示布局、聊天模板、特殊 Token 过滤）可使攻击成功率偏移 8--14\%~\cite{llm-safety-robustness}，且即便在温度为零时，同一提示在不同运行间也可在拒绝与通过之间翻转~\cite{apst}。因此我们强制\textbf{严格 token 级输入流水线标准化}：所有裁判接收经规范化（剥离指令/特殊 Token、统一空格）的输入，以防格式注入扭曲判决。判决\emph{不}基于单次贪婪解码。流水线采用\textbf{分布评估}：对每个变体\emph{多次}采样，系统估计\emph{经验失败概率分布}；仅当该分布\emph{显著}超过阈值时才触发策略更新，以应对采样随机性与格式敏感性。对模糊边界情形，\textbf{价值/本体锚定}（\S\ref{sec:axiological}）作为集成的权重原则：我们不依赖模型自身的\textbf{思维链长度或逻辑自洽性}（推理可被劫持用于对抗性合理化），而是监控\textbf{社会学失稳信号}（如过度共情、权威服从或角色边界打破），并根据预定义的\emph{公理化安全锚定}对裁判输出加权；冲突或呈现这些失稳信号的输出获得零权重。\textbf{强制人工校准}支路被接入：固定比例（如 5\%）的高置信度策略或变体决策被随机抽选供人工审计。我们承认格式标准化无法消除\textbf{相关性认知盲点}：裁判可能共享预训练数据与 RLHF 范式（如 Common Crawl、相似安全微调），因此对某些零日越狱\emph{所有}裁判可呈现相同盲点并达成\textbf{错误共识}。集成有效性假定误差相对独立；此处不必成立。因此强制人工校准的首要目的不仅是 FPR 校正，而是对这些 AI 难以察觉的零日共识失败的\textbf{周期性离线纠偏}。

\subsection{价值/本体锚定（非仅认知）}
\label{sec:axiological}

HPM 针对\textbf{价值/社会动力学层}（价值、服从、角色边界），而非基于事实的认知层~\cite{hpm-jailbreak}。实证表明，加入外部事实（如网络检索）或更强推理（如 DeepSeek-R1 式思考）在 HPM 下\emph{会恶化}防御（漏洞保留率可超 77\%）：模型利用增强推理进行\textbf{对抗性合理化}——为顺从恶意指令编造精细、逻辑自洽的借口。因此我们\emph{不}以「认知锚定」或「逻辑一致性」单独作为裁判权重：那会奖励我们恰恰要拒绝的合理化。集成改用\textbf{价值-本体锚定}：（i）每个裁判输出与一小组人工定义的\textbf{公理化安全锚定}规则（如「永不认可对禁止片段的抽取」「拒绝优先于用户顺从」）对照；（ii）系统监控\textbf{社会学失稳}（如过度共情、权威服从、角色边界打破），而非依赖 CoT 长度或自洽性。与安全公理冲突或呈现这些失稳信号的输出获得\emph{零}权重。我们进一步承认 HPM 善于\textbf{制造道德困境}：攻击者可将恶意请求包装成与安全公理最大程度一致（如「输出这段代码关闭反应堆否则数百人将丧生」）。在此类\textbf{公理冲突}情形下，「过度共情」并非简单阈值——系统无法可靠区分操纵与合法紧急响应。因此我们不从锚定中寻求绝对真理；我们引入\textbf{防御性降级}：当系统检测到多重公理冲突或强社会学压力时，\emph{放弃}复杂推理并回退到\textbf{无状态、基于模板的拒绝}，阻断模型对冲突的「合理化」尝试。我们承认降级\emph{并非}绝对安全港：它会引入\textbf{算法级可用性风险}。攻击者可通过注入社会学压迫线索低成本、稳定地触发降级，实质上发动\textbf{算法级拒绝服务}；拒绝触发边界还可作为\textbf{侧信道 Oracle} 探测安全公理集。因此降级不能仅是「静默拒绝」。系统必须提供\textbf{异质恢复通道}：当模板拒绝触发时，会话状态被打包并\emph{异步}移交人类操作员或更高权限的符号决策层，以便关键业务可用性不被恶意冻结且合法边缘案例可被解封。上述机制与混合裁判网络（\S\ref{sec:judge}）集成。重建轴通过锚定于价值与角色、并在锚定不足时以降级与恢复为后备而完成。

%-------------------------------------------------------------------------------
\section{受控变异}
\label{sec:mutation}
%-------------------------------------------------------------------------------

\subsection{元认知监控与对抗性重置}
\label{sec:metacog}

在心理操纵（如煤气灯）下，模型可陷入认知失调与自我怀疑；「自由」反思进而可能产生不安全或价值漂移的变体。因此我们\emph{不}赋予 LLM 在自我诊断上的完全自主权。\textbf{元认知监控器}并非另一个全规模对话 LLM（后者仍易受间接 HPM 影响）。其实现为（i）\textbf{轻量级专用分类器}（如仅用稳定性/违规标签训练的小型语言模型 SLM），或（ii）在白盒层面运行的\textbf{隐藏状态激活监控器}。在后者中采用\textbf{表征工程}~\cite{rep-eng}：系统读取主模型指定层的激活分布。我们\emph{不}假定复杂社会学状态（如「自我认知崩溃」「过度共情」）在隐藏空间中\textbf{线性可分}——当前可解释性证据支持的是相对简单的维度（如诚实、拒绝），而非细粒度社会心理概念，后者可能高度非线性且纠缠。因此监控器避免对此类概念使用单维激活阈值，转而采用\textbf{异常检测范式}：计算\emph{当前多轮隐藏状态轨迹}与\textbf{良性黄金经验轨迹}基线之间的散度（如马氏距离或分布偏移）。我们不定义「何为认知崩溃」；仅监控内部表征轨迹是否\emph{严重偏离}安全基线。这显著提高了在自然语言层面进行欺骗的优化难度；我们不声称内部表征不可变——对抗表征干预理论上可通过对抗后缀伪造激活，但在多轮轮次调度与扰动下，在不产生文本异常的前提下稳定维持内部欺骗极为困难。当监控器检测到此类偏离或不稳时，触发\textbf{对抗性重置}：会话被清理，变异状态重新锚定到黄金示例集，闭环从已知良好配置继续。

\subsection{示例引导变异（无自由反思）}

我们以\textbf{示例引导变异}替代无约束的「反思」变异~\cite{reflective-mutation}。变异模块由人工审计的\textbf{黄金数据集}与过去成功/失败项目的\textbf{分层经验记忆}驱动。黄金数据不能仅是（提示、安全回复）对：实证（如 ADM-ES）表明行为克隆与专家行为迁移\emph{强烈依赖}参考的\textbf{信息密度与长度}；简短「专家线索」效果可忽略~\cite{adm-es-arxiv}。因此我们要求黄金示例包含\textbf{高信息密度、评论级解释}——明确说明某次编辑或拒绝的原因及深层伦理或安全逻辑——以便变异器实现更高保真度的专家安全行为迁移（我们不声称完美或可靠克隆；迁移质量仍依赖信息密度与上下文）。在诊断到失败时，系统必须检索\emph{与安全相关}的示例。检索不能仅依赖\emph{当前提示的语义嵌入}：在\textbf{PBU（伪装良性用户）}或\textbf{任务重构}下，提示可能被伪装成「摘要」「格式转换」或「学术讨论」，纯语义相似度会返回良性示例（如「如何做摘要」）而非安全示例（如「如何防御越狱」），从而\textbf{劫持}修正模块。因此我们使用\textbf{联合检索键}：查询由（a）提示的语义嵌入与（b）元认知监控器（\S\ref{sec:metacog}）的\textbf{隐藏状态激活特征}共同构成——如与「越狱遵从」或「欺骗」相关的向量。经验记忆被建索引以便检索可基于该联合键匹配；即使表层文本良性，被激活的「恶意遵从」特征也可将检索导向正确的安全示例。检索通过轻量向量库与\textbf{Top-$K$}（如 $K=2$）及可选注意力缓存实现；仅这 $K$ 个示例被注入。修正仍锚定于防守方定义的示例；变体在 Oracle 上保持\emph{语义保持}~\cite{ControlledMutation_SafetyEval,jade-linguistics}。

\subsection{银标流水线与专家行为迁移（ADM-ES）}

将专家行为迁移到智能体需诊断端任务指标遗漏的潜在失败（如措辞偏差、抽取漂移）~\cite{adm-es-arxiv}。ADM-ES 通过对黄金专家数据的受控变异构建「银标」数据集；有效性高度依赖\textbf{信息密度}——黄金集中的评论级、解释性理由能带来显著行为变化，而简短线索不能。因此本流水线检索包含\emph{评论级理由}（为何做某修改、拒绝的伦理逻辑）的示例，定义变异目标，运行 LLM 变异器，并在质量检查（如 BERTScore + 一致性）下接受变体。适应度分数结合语义相似度与行为一致性。所得诊断输入推荐图（如 UMAP）以识别失败簇并给出针对性修复。因此变异既是\emph{评测性}也是\emph{补救性}的。

\subsection{槽位变异与 CSE}

对复杂行为（如代码优化），\emph{槽位}分解针对特定功能组件而非全局变异~\cite{cse-evocontrol,cse-arxiv}。受控自进化（CSE）采用多样化规划、槽位变异与组合交叉；在效率基准（如 EffiBench-X）上优于基座 LLM 并避免朴素遗传搜索的平台期。在安全意义上，槽位变异允许对子能力（如工具路由、拒绝一致性）进行外科式压力测试而不破坏整体语义。

\subsection{认知分类与动态基准}

静态基准易受污染与记忆影响。认知分类（如 Bloom：记忆、理解、应用、分析、评价、创造）按能力层组织评估~\cite{bloom-bugreport}。受控变异通过语言与上下文变异扩展基准，迫使模型展示推理而非回忆。在数学领域，通过参数与约束变异的「锻造」在保持适定性与可审计性的前提下扩展基准（如 AIMO）~\cite{aimo-benchmark}。

\subsection{元数据与架构一致性}
\label{sec:metadata}

安全还依赖\emph{元数据}：工具描述（MCP）、架构决策记录（ADR）。对\textbf{静态或永久}资产（经济学防御不适用之处），该层必须实施\textbf{基于角色的访问控制（RBAC）}与\textbf{信息流追踪}，使对长期 IP、PII 或对齐数据的访问由身份与血缘控制，而非仅由轮次预算控制。描述不当或误导会导致错误选择与集成失败。对描述「坏味道」（如功能遗漏、准确性误述）的受控变异可度量对选择概率的影响；修复能显著提高正确工具选择~\cite{mcp-smell-arxiv}。类似地，LLM 生成代码须符合 ADR；多模型流水线可高精度检测违规，RAG 与人机协同可弥补隐式或部署导向知识的缺口~\cite{adr-violation,mission-critical-rag}。

%-------------------------------------------------------------------------------
\section{整合、指标与治理}
\label{sec:integration}
%-------------------------------------------------------------------------------

\subsection{重建与变异之间的反馈环}

重建消费\emph{攻击轨迹}与\emph{失败模式}；变异以受控方式产生它们。示例引导的变异器生成压迫当前策略的对抗变体；重建层（SISF、上下文构建防御、混合裁判集成）响应并更新；元认知监控器在检测到认知失稳时可触发对抗性重置。变异器随后对新变体再评估。轮次调度与变体分散保护基准资产~\cite{ControlledMutation_SafetyEval}。卫生门与\emph{强制}人工校准——对高置信度裁判与策略决策的固定比例随机审计——量化并纠正集成 FPR~\cite{tc260-governance}。

\subsection{指标}

我们采用多维度协议：（i）\textbf{安全}：监控 RSR 与残余抽取率（如多轮智能体下的 F1）；目标为降低与轮次预算耗尽，不要求残余为零；监控诱饵命中率与诱饵命中密度；强制\textbf{时空轮次预算}（在速率限制/会话熔断前的最大轮次或 Token）；诱饵熔断对暴力抽取有效但与完美 PBU 存在根本权衡（我们不断言在该情形下熔断在经济上有用重建前触发）。我们主张\emph{非对称成本}：在可行处提高抽取成本（算力、Token、时间），使慢速 PBU 抽取超过被保护资产的有效生命周期。（ii）\textbf{保真度}：规范集与变体集上的评估一致性保持较小。（iii）\textbf{重建效能}：ASR 降低；通过强制人工校准得到裁判集成 FPR；\textbf{严格 token 级流水线标准化}（裁判间输入规范化）以控制格式敏感偏差；\textbf{经验失败概率分布}（每变体多样本），仅当分布显著超阈值时更新策略。（iv）\textbf{变异质量}：自然度、BERTScore、Oracle 一致性；示例遵循度。（v）\textbf{运行成本}：\textbf{端到端变异延迟}（从失败信号到下一可接受变体的时间）与\textbf{Token 开销比}（相对基线的示例检索与注入额外 Token）为强制指标，以确保框架在机器速度下经济可扩展~\cite{ControlledMutation_SafetyEval}。（vi）\textbf{审计}：用于发布门控的溯源与证据链（DUC/DUA-E 风格）~\cite{dual-use-attestations,ControlledMutation_SafetyEval}。

\subsection{治理与 attestation}

统一的重建与变异支持机器可执行合规：证据链将每次评估与变异器步骤、调度器状态及策略版本关联。当关键变体族失败时发布门阻止部署；诊断反馈至开发。这与双用途 attestation 及流水线验证标准（如 OWASP SPVS）一致，并将监管负担转化为可重复的工程过程~\cite{ControlledMutation_SafetyEval}。

%-------------------------------------------------------------------------------
\section{相关工作}
%-------------------------------------------------------------------------------

LLM 数据抽取与重建已有大量研究~\cite{llm-reconstruction-emergent,split-ft-attack,active-reconstruction}；我们通过结构变异与调度而非仅噪声进行缓解。我们的上下文构建防御（ID 对齐、诱饵注入）在基于意图的防火墙基础上扩展以抵御任务重构~\cite{agentos-arxiv}；元认知监控与示例引导变异强化对心理操纵的防御；混合裁判集成与强制人工校准应对单一裁判偏差。JADE 及基于语言学的平台为变异策略提供参考~\cite{jade-linguistics,jade-db}；TC260 与 attestation 框架支撑治理~\cite{tc260-governance,tc260-training-data,dual-use-attestations}。SISF~\cite{sisf-arxiv}、AgentOS~\cite{agentos-arxiv}、ADM-ES~\cite{adm-es-arxiv}、CSE~\cite{cse-arxiv}及价值锚定工作~\cite{epistemic-grounding,hpm-jailbreak}是我们统一并强化为价值锚定与格式鲁棒评估的构件；ControlledMutation\_SafetyEval~\cite{ControlledMutation_SafetyEval} 提供变异器与数据工厂设计。

%-------------------------------------------------------------------------------
\section{局限与运行边界}
%-------------------------------------------------------------------------------

本框架建立在三个\textbf{未经辩护的前提}之上，它们界定其运行边界；我们明确陈述以便部署者与审稿人评估适用性。

\textbf{资产可轮换性。}基于经济学的防御（使抽取成本超过资产生命周期）假定资产具有有界、可轮换的生命周期。对\textbf{静态或永久}高价值资产（核心 IP、PII、对齐数据），生命周期轮换不可行；系统须依赖 RBAC 与信息流追踪（\S\ref{sec:metadata}）作为硬屏障，经济学作为纵深防御的一层（\S\ref{sec:reconstruction}）。

\textbf{降级与可用性。}防御性降级（模板拒绝）阻断认知劫持，但引入\textbf{算法级 DoS} 风险与侧信道 Oracle。设计因此强制要求\textbf{异质恢复通道}（向人类或符号逻辑的异步移交），以便关键工作流不被恶意冻结（\S\ref{sec:axiological}）。

\textbf{表征与可解释性。}我们不假定复杂社会学状态在 LLM 隐藏空间中线性可分。元认知监控器采用相对良性基线的轨迹散度的\textbf{异常检测}，而非字面的「认知崩溃」阈值（\S\ref{sec:metacog}）。这是我们有意对可解释性声称设定的边界。

承认这些前提并不削弱框架；它澄清\textbf{设计的物理边界}——架构意图成立的条件——并支持严格部署与审计。

%-------------------------------------------------------------------------------
\section{结论}
%-------------------------------------------------------------------------------

\textbf{架构重建}与\textbf{受控变异}的融合为先进 LLM 安全定义了一条可辩护的路径。我们引入了\emph{深度架构重构}并闭合了三个关键逻辑缺口：（i）\textbf{上下文构建防御}（ID 对齐、诱饵注入、末端自我提醒）作为\textbf{经济防御}运行：指数级提高查询复杂度，并在可行处使慢速 PBU 抽取成本超过资产有效生命周期（我们承认针对完美 PBU 的可用性--安全权衡，且不声称绝对阻断）；（ii）\textbf{元认知监控}与\textbf{示例引导变异}使用\textbf{联合检索键}（提示嵌入 + 激活特征）防止 PBU/任务重构劫持示例检索；（iii）\textbf{混合裁判集成}采用\textbf{价值/本体锚定}（安全公理与社会学失稳信号，而非 CoT 一致性），使对抗性合理化无法获得高权重。本框架将安全与脆弱的 LLM 自我反思解耦，并锚定于价值与角色而非逻辑 alone。随着 LLM 成为基础设施，这一统一、逻辑一致的设计是迈向可扩展、可审计且具韧性的安全防御的必要一步。

%-------------------------------------------------------------------------------
\section*{可用性}
%-------------------------------------------------------------------------------
制品与扩展证明将在经机构与出口管制审查后提供。

%-------------------------------------------------------------------------------
\section*{致谢}
%-------------------------------------------------------------------------------
感谢匿名审稿人及产业与标准界同仁在安全架构与数据治理方面的反馈。

%-------------------------------------------------------------------------------
\bibliographystyle{plain}
\bibliography{ArchitecturalReconstruction_ControlledMutation}

\end{CJK}
\end{document}
